% ============================================================================
%  chapter14.tex
%  Text Parsing and State Machines
% ============================================================================

\chapter{Text Parsing and State Machines}

\section{Why Parsing Is a Systems Skill}

Every program that reads text must \emph{parse} it: take raw bytes and
extract the structure the text encodes. Configuration files, HTTP
headers, JSON, CSV, command lines, and every wire protocol are all
text with rules, and the code that reads them is a \emph{parser}.
Parsers are everywhere, they are written by hand more often than you
would expect, and they are the classic home of two bugs: crashes on
malformed input, and silent misinterpretation of ``almost valid''
input.

This chapter builds a real parser end to end. The example is the
command language of \emph{MiniDB}, a small key-value database written
in C that lives beside this book (in the workspace's
\code{MiniDB-deepseek-implementation} directory). MiniDB accepts
commands like \code{INSERT 42 100} and \code{SELECT 7} from a prompt;
before any data is touched, every line must be \emph{lexed} into
tokens and \emph{parsed} into a command. We will study its lexer and
parser as the running example, and adapt its code to this book's
style.

\begin{verbatim}
   The layers of parsing:

     raw text:   "INSERT 42 100\n"
                     |
                     v   LEX (tokenize, one char at a time)
     tokens:     (KEYWORD: INSERT) (INTEGER: 42) (INTEGER: 100) (END)
                     |
                     v   PARSE (validate shape, fill the struct)
     command:    { type = COMMAND_INSERT, key = 42, value = 100 }

   - A TOKEN is the smallest meaningful unit (a word, a number, a
     symbol). The LEXER (tokenizer) splits text into tokens and
     discards what is insignificant (whitespace).
   - The PARSER checks the token stream against the grammar and fills
     a struct (here: a Command). The database then executes it.
\end{verbatim}

\begin{osconcept}
Parsing is a systems skill because text is the universal interface
between programs: every configuration file, log line, and protocol
message your daemon reads arrives as bytes and must be understood
safely. The OS delivers you bytes; it will not protect you from
malformed input. Turning unreliable text into reliable structure, and
failing cleanly when the text is wrong, is the programmer's own
problem---which is exactly why parsers are where so much care lives.
\end{osconcept}

\section{Tokens and the Token Stream}

A \emph{token} is a run of characters that counts as one unit: an
identifier, a number, a quoted string, a symbol. The tokenizer's only
job is splitting; deciding what tokens \emph{mean} belongs to the
parser. MiniDB has just three token kinds, and that simplicity is
deliberate:

\begin{verbatim}
   MiniDB's language, in one screen:

     INSERT <key> <value>     store a key/value pair
     SELECT <key>             fetch one key
     DELETE <key>             remove one key
     RANGE  <lo> <hi>         fetch all keys in [lo, hi]
     PRINT                    dump the whole database
     EXIT                     leave the REPL
     HELP                     list the commands

   Every command is one verb plus a FIXED number of integers.
   That tiny grammar keeps the parser flat (no tree needed yet).
\end{verbatim}

\begin{lstlisting}[caption={The token types and the token struct.}]
typedef enum {
    TOKEN_KEYWORD,     /* names a command verb; read "keyword" */
    TOKEN_INTEGER,     /* a bare decimal literal; read "int_value" */
    TOKEN_END          /* sentinel: the stream ends here        */
} TokenType;

typedef enum {
    KEYWORD_INSERT, KEYWORD_SELECT, KEYWORD_DELETE, KEYWORD_RANGE,
    KEYWORD_PRINT, KEYWORD_EXIT, KEYWORD_HELP, KEYWORD_UNKNOWN
} Keyword;

typedef struct {
    TokenType type;
    Keyword   keyword;    /* valid when type == TOKEN_KEYWORD */
    int       int_value;  /* valid when type == TOKEN_INTEGER */
} Token;
\end{lstlisting}

The field comments make the \emph{discriminator} explicit: which member
of the struct is meaningful depends on \code{type}. That is the tagged
union pattern of Chapter~3, back for its first serious workout.

Two design decisions in these few lines will matter through the whole
chapter. First, every stream \emph{ends with a sentinel}
(\code{TOKEN\_END}): the parser can check for ``the stream is done''
without knowing the length by other means, and can never read past the
last real token. Second, the lexer \emph{allocates nothing}: the
caller owns the token array, and the lexer just fills it up to a fixed
bound. A command line has at most a handful of tokens, so a fixed
array is not a limitation---it is a simplification that removes an
entire class of memory bugs from the parser.

\begin{verbatim}
   "  exit  "  ====>  (KEYWORD: EXIT) (END)      whitespace vanishes
   "INSERT 42" ====>  (KEYWORD: INSERT) (INTEGER: 42) (END)
   "select 7"  ====>  (KEYWORD: SELECT) (INTEGER: 7)  (END)
   (keywords are case-insensitive; MINIDB matches them uppercased)
\end{verbatim}

\section{The Lexer as a State Machine}

The lexer is a tiny character-level \emph{state machine}: it walks the
input once, left to right, and each character moves it into exactly one
state. A \emph{state machine} (in its simplest form, a \emph{finite
automaton}) has a finite set of states, one start state, and
\emph{transitions}---rules saying what happens when a character
arrives in each state. Some states are \emph{accepting}: if the input
ends there, the input was valid.

\begin{verbatim}
   MiniDB's lexer, drawn as a state machine:

                              isspace (skip, repeat)
                              |
        START --------------->+<-- SKIP_WS
         |
         | isdigit                     | any non-whitespace
         v                             v
       INTEGER -----------------> (back to START for the next char)
         |  overflow                       ^
         v                                 |
       ERROR <--isalpha? no, no match-- WORD (matched a keyword -> token)
         |
         +-- "any character outside the grammar" also lands here

   The accepting state is END OF INPUT: when the scan runs out of
   characters, it appends TOKEN_END. Everything else either emits a
   token and loops, or enters ERROR and aborts the whole scan.

   One character, one decision, never a lookahead, never a backup:
   that is exactly what a finite automaton is.
\end{verbatim}

The lexer below is the MiniDB lexer, faithful in structure and
condensed in one place: the original spells out each of the seven
keywords with its own comparison block for teaching clarity; here we
match against a table in a loop. The states---skip whitespace,
integer, word, stray symbol, end---are unchanged.

\begin{lstlisting}[caption={The MiniDB lexer (condensed from lexer.c).}]
#include <ctype.h>
#include <limits.h>
#include <string.h>

#define LEXER_MAX_TOKENS 16

typedef enum {
    LEX_OK,
    LEX_ERROR_UNKNOWN_WORD,     /* unrecognized word or stray symbol  */
    LEX_ERROR_BAD_NUMBER,       /* digit run that would overflow int  */
    LEX_ERROR_TOO_MANY_TOKENS   /* more tokens than the array holds   */
} LexResult;

static const char *const KEYWORD_TABLE[KEYWORD_UNKNOWN] = {
    "INSERT", "SELECT", "DELETE", "RANGE",
    "PRINT", "EXIT", "HELP"
};

/* The ONLY place tokens enter the output array. Bounds-checked. */
static LexResult append_token(Token *tokens, size_t *token_count, Token t) {
    if (*token_count >= LEXER_MAX_TOKENS) {
        return LEX_ERROR_TOO_MANY_TOKENS;
    }
    tokens[*token_count] = t;
    *token_count = *token_count + 1;
    return LEX_OK;
}

LexResult lex(const char *input, Token *tokens, size_t *token_count) {
    const char *p;
    size_t count;

    p = input;
    count = 0;

    while (*p != '\0') {

        /* ---- State: skip whitespace (never significant) ---- */
        if (isspace((unsigned char) *p)) {
            p = p + 1;
            continue;
        }

        /* ---- State: integer literal ---- */
        if (isdigit((unsigned char) *p)) {
            Token t;
            int value;

            value = 0;
            while (isdigit((unsigned char) *p)) {
                int digit;

                digit = *p - '0';
                /* Overflow guard: stop before value * 10 overflows. */
                if (value > (INT_MAX - digit) / 10) {
                    return LEX_ERROR_BAD_NUMBER;
                }
                value = value * 10 + digit;
                p = p + 1;
            }

            t.type = TOKEN_INTEGER;
            t.int_value = value;
            if (append_token(tokens, &count, t) != LEX_OK) {
                return LEX_ERROR_TOO_MANY_TOKENS;
            }
            continue;
        }

        /* ---- State: word (a keyword, upcased for matching) ---- */
        if (isalpha((unsigned char) *p)) {
            char word[8];     /* longest keyword + NUL */
            size_t wlen;
            int matched;
            size_t i;

            wlen = 0;
            while (isalpha((unsigned char) *p)) {
                if (wlen >= sizeof(word) - 1) {
                    return LEX_ERROR_UNKNOWN_WORD;   /* too long */
                }
                word[wlen] = (char) toupper((unsigned char) *p);
                wlen = wlen + 1;
                p = p + 1;
            }
            word[wlen] = '\0';

            matched = 0;
            for (i = 0; i < (size_t) KEYWORD_UNKNOWN; i++) {
                if (wlen == strlen(KEYWORD_TABLE[i]) &&
                    memcmp(word, KEYWORD_TABLE[i], wlen) == 0) {
                    Token t;

                    t.type = TOKEN_KEYWORD;
                    t.keyword = (Keyword) i;
                    if (append_token(tokens, &count, t) != LEX_OK) {
                        return LEX_ERROR_TOO_MANY_TOKENS;
                    }
                    matched = 1;
                    break;
                }
            }
            if (matched == 0) {
                return LEX_ERROR_UNKNOWN_WORD;
            }
            continue;
        }

        /* ---- State: stray symbol (not part of the language) ---- */
        return LEX_ERROR_UNKNOWN_WORD;
    }

    /* ---- State: end of input -- emit the sentinel ---- */
    {
        Token end;

        end.type = TOKEN_END;
        end.keyword = KEYWORD_UNKNOWN;
        end.int_value = 0;
        if (append_token(tokens, &count, end) != LEX_OK) {
            return LEX_ERROR_TOO_MANY_TOKENS;
        }
    }

    *token_count = count;
    return LEX_OK;
}
\end{lstlisting}

Read the code against the diagram and three facts become visible. The
overflow guard is a genuine state-machine error transition: it aborts
mid-number rather than emit a wrong token. The word buffer is bounded
(\code{word[8]}, and the capacity check before every write), so a
100-character ``word'' is rejected, never overflowed. And on \emph{any}
error the function returns immediately; the caller must treat the
partially-filled stream as garbage. There is no recovery-and-continue;
MiniDB chooses fail-fast, which is the right choice for a command
line.

\begin{verbatim}
   The states, recapped in one line each:

     SKIP_WS    whitespace is never significant
     INTEGER    consume digits, guard overflow, emit INTEGER
     WORD       upcase and match against the keyword table
     STRAY      any other byte is a one-character unknown word
     END        append the TOKEN_END sentinel and finish

   No dynamic allocation anywhere. The caller owns the array; the
   lexer owns nothing. That one decision is why the whole parser can
   be reasoned about with plain C arrays and explicit bounds.
\end{verbatim}

\begin{ruleofthumb}[Draw the machine before writing the loop]
Any parser with more than one kind of token deserves a state diagram
first. Sketch the states, walk the happy path, walk two malformed
inputs, and only then write the \code{while}. The diagram is the
specification; the code is the transcription. Parsers written the
other way around are how ``almost valid'' input silently becomes
wrong data.
\end{ruleofthumb}

\section{The Parser: Validating the Token Stream}

Where the lexer is a \emph{scanner}, the parser is a \emph{validator}.
It answers two questions: \emph{is this stream a legal command, and
what does it mean?} It does not look at characters again---the lexer
already did that work. It walks the token array with a cursor and
checks the stream's \emph{shape} and \emph{arity} against the grammar.

MiniDB's grammar is the smallest possible: one verb, plus a fixed
count of integer arguments. There is nothing to nest, so the parser
barely needs recursion at all. It works in three moves:

\begin{verbatim}
   The parser's three moves:

   1. DISPATCH on tokens[0].  It must be a KEYWORD; a switch on its
      value selects the rule. A bare integer first is not a verb:
      PARSE_ERROR_UNKNOWN_COMMAND.

   2. EXACT-ARITY CHECK.  Each command calls parse_fixed_arguments()
      with the number of integers it expects. One helper consumes
      exactly that many integers, then REQUIRES the very next token
      to be END. Three failures in one place:
         - stream ends early          -> MISSING_ARGUMENT
         - a non-integer where an int belongs -> BAD_ARGUMENT
         - anything left over         -> EXTRA_ARGUMENT

   3. BUILD THE Command.  On success, the integers are copied into the
      right fields and command->type is set.
\end{verbatim}

\begin{lstlisting}[caption={The exact-arity engine shared by every command.}]
typedef enum {
    PARSE_OK,
    PARSE_ERROR_EMPTY,
    PARSE_ERROR_UNKNOWN_COMMAND,
    PARSE_ERROR_MISSING_ARGUMENT,
    PARSE_ERROR_EXTRA_ARGUMENT,
    PARSE_ERROR_BAD_ARGUMENT
} ParseResult;

/* Consume exactly argc integers starting at tokens[start]; then the
   next token MUST be TOKEN_END. Results land in args[].          */
static ParseResult parse_fixed_arguments(const Token *tokens,
                                         size_t token_count,
                                         size_t start, size_t argc,
                                         int *args) {
    size_t i;

    for (i = 0; i < argc; i++) {
        if (start >= token_count) {
            return PARSE_ERROR_MISSING_ARGUMENT;   /* no sentinel at all */
        }
        if (tokens[start].type == TOKEN_END) {
            return PARSE_ERROR_MISSING_ARGUMENT;   /* stream ran short */
        }
        if (tokens[start].type != TOKEN_INTEGER) {
            return PARSE_ERROR_BAD_ARGUMENT;       /* wrong kind of token */
        }
        args[i] = tokens[start].int_value;
        start = start + 1;
    }

    if (start >= token_count) {
        return PARSE_ERROR_MISSING_ARGUMENT;
    }
    if (tokens[start].type != TOKEN_END) {
        return PARSE_ERROR_EXTRA_ARGUMENT;         /* too many arguments */
    }
    return PARSE_OK;
}
\end{lstlisting}

\begin{verbatim}
   The same walker, three failing inputs:

     INSERT 42          ->  expects 2 ints, finds 1, then END
                            -> MISSING_ARGUMENT
     INSERT 42 100 5    ->  two ints consumed, then an INTEGER
                            where END must be -> EXTRA_ARGUMENT
     INSERT x 100       ->  an integer demanded, KEYWORD found
                            -> BAD_ARGUMENT

   One helper, three errors: the arity is checked by construction,
   because the "then it must end" line is part of the same function.
\end{verbatim}

The dispatch is a plain switch over the keyword, each case stating its
arity in one line---that number \emph{is} the grammar rule:

\begin{lstlisting}[caption={parse\_command: dispatch on the keyword, fill the Command.}]
typedef enum {
    COMMAND_INSERT, COMMAND_SELECT, COMMAND_DELETE, COMMAND_RANGE,
    COMMAND_PRINT, COMMAND_EXIT, COMMAND_HELP, COMMAND_NONE
} CommandType;

/* The proto-AST: one flat node, because every statement is one verb
   plus at most two integers. */
typedef struct {
    CommandType type;
    int key;      /* INSERT/SELECT/DELETE/RANGE: the primary key   */
    int value;    /* INSERT only: the value to store               */
    int key_end;  /* RANGE only: the inclusive upper bound         */
} Command;

ParseResult parse_command(const Token *tokens, size_t token_count,
                          Command *command) {
    int args[2];               /* largest arity in the grammar is two */
    ParseResult r;

    /* Initialize every field first: unused fields stay defined. */
    command->type = COMMAND_NONE;
    command->key = 0;
    command->value = 0;
    command->key_end = 0;

    if (token_count == 0 || tokens[0].type == TOKEN_END) {
        return PARSE_ERROR_EMPTY;          /* nothing there at all */
    }
    if (tokens[0].type != TOKEN_KEYWORD) {
        return PARSE_ERROR_UNKNOWN_COMMAND;/* an integer is not a verb */
    }

    switch (tokens[0].keyword) {
    case KEYWORD_INSERT:
        r = parse_fixed_arguments(tokens, token_count, 1, 2, args);
        if (r != PARSE_OK) return r;
        command->type = COMMAND_INSERT;
        command->key = args[0];
        command->value = args[1];
        return PARSE_OK;

    case KEYWORD_SELECT:
        r = parse_fixed_arguments(tokens, token_count, 1, 1, args);
        if (r != PARSE_OK) return r;
        command->type = COMMAND_SELECT;
        command->key = args[0];
        return PARSE_OK;

    case KEYWORD_DELETE:
        r = parse_fixed_arguments(tokens, token_count, 1, 1, args);
        if (r != PARSE_OK) return r;
        command->type = COMMAND_DELETE;
        command->key = args[0];
        return PARSE_OK;

    case KEYWORD_RANGE:
        r = parse_fixed_arguments(tokens, token_count, 1, 2, args);
        if (r != PARSE_OK) return r;
        command->type = COMMAND_RANGE;
        command->key = args[0];
        command->key_end = args[1];
        return PARSE_OK;

    /* The three zero-argument commands all reduce to: END must follow. */
    case KEYWORD_PRINT:
    case KEYWORD_EXIT:
    case KEYWORD_HELP:
        r = parse_fixed_arguments(tokens, token_count, 1, 0, args);
        if (r != PARSE_OK) return r;
        command->type = (CommandType) (COMMAND_PRINT +
            (tokens[0].keyword - KEYWORD_PRINT));
        return PARSE_OK;

    default:
        return PARSE_ERROR_UNKNOWN_COMMAND;
    }
}
\end{lstlisting}

\begin{verbatim}
   The token-to-Command mapping, drawn:

     tokens:  (KEYWORD: INSERT) (INTEGER: 42) (INTEGER: 100) (END)
                                     |              |
                                     v              v
     Command: { type = COMMAND_INSERT, key = 42, value = 100, key_end = 0 }

     tokens:  (KEYWORD: RANGE) (INTEGER: 10) (INTEGER: 20) (END)
                                     |              |
                                     v              v
     Command: { type = COMMAND_RANGE, key = 10, value = 0, key_end = 20 }
\end{verbatim}

Calling the output a \emph{proto-}AST is a hint about the future. A
real query such as \code{SELECT price FROM items WHERE id = 42} would
need nested expressions, so \code{Command} would grow pointers to
sub-structures and the parser would recurse. MiniDB's statements are
all one verb plus at most two integers, so the ``tree'' collapses to
one node---but the seam where tree-ness would return is identified,
and that is exactly the place our Chapter~4 trees would slot in.

\section{The REPL: Read, Gate, Execute}

\subsection{What a REPL Is}

A \emph{REPL}---\textbf{R}ead-\textbf{E}val-\textbf{P}rint
\textbf{L}oop---is an interactive top-level loop. It is the oldest user
interface in programming, invented for \emph{Lisp} and later adopted
by early BASIC and every language console since (Python, Ruby,
JavaScript, SQL: all of them drop you into a REPL). It means exactly
what its name says, three phases repeated forever:

\begin{verbatim}
   the three phases of every REPL:

     READ    get one complete unit of input   (one line, here)
     EVAL    do the work on it                (lex + parse + execute)
     PRINT   show the outcome to the user     (the reply message)
                ^
                +---- and loop, forever ----------+

   Contrast with a BATCH program, which reads all input, computes, and
   prints once: a REPL is *interactive* -- it must answer each input
   before waiting for the next. That single requirement shapes
   everything: prompt, flush, per-line errors, loop.
\end{verbatim}

\textbf{What Lisp was.} Lisp (for \emph{LISt Processing}) is a family
of programming languages designed by John McCarthy starting in 1958---
the second-oldest high-level language after Fortran. Its founding idea
was that the \emph{list} is the universal structure: a Lisp program is
itself a list of symbols, so a program can read another program, alter
it, and run the result. That property (\emph{homoiconicity}, ``code is
data'') is why Lisp became the language of artificial-intelligence
research for three decades. Lisp also gave computing two gifts that
every language now takes for granted: \emph{garbage collection}
(automatic reclamation of memory, which we build pieces of in
Chapter~13) and the \emph{REPL} itself. A Lisp programmer did not
write, compile, and run---they talked to the machine, one expression
at a time, and the machine answered immediately. Interactive
development was Lisp's invention, and the MiniDB console is its
descendant.

Two consequences follow from ``interactive.'' First, the loop must
\emph{survive} bad input: one bad line cannot kill the session, so
every error becomes ``print a message, continue.'' Second, the loop
must announce itself with a \emph{prompt} (\code{minidb> }), and that
prompt must be \emph{flushed}: when stdout is a pipe rather than a
terminal, it is fully buffered, and without \code{fflush(stdout)} the
prompt would sit in a buffer while the program waited for input it
never visibly requested. Interactive software must flush.

\begin{verbatim}
   The REPL is the front end of every interactive system -- a shell,
   a database console, a debugger, a calculator. MiniDB's is small,
   but it is the same shape as all of them.
\end{verbatim}

\subsection{REPL Gating: Every Stage Guards the Next}

Now the word \emph{gating}. The REPL we are about to read does not
pipe raw input to the database. It passes each line through a sequence
of validation stages, and each stage is a \emph{gate}: it either
converts the input into a more structured form and passes it through,
or it rejects the line and the whole pipeline stops there. The
database never sees a bad line, because a bad line never gets through
two gates to reach it.

\begin{verbatim}
   Gating, drawn as a funnel:

      raw line
        |
        v
   +----------+   pass?   +----------+   pass?   +-----------+
   | LEX gate |---------->| PARSE    |---------->| EXECUTE   |
   | chars -> | fail:     | tokens ->| fail:     | command - |
   | tokens   | print,    | Command  | print,    | > storage |
   +----------+ continue  +----------+ continue  +-----------+
      |                       |                       |
      v                       v                       v
   "Lex error."           "Parse error."           the reply

   Two walls protect the storage engine from the user: a line must
   survive the lexer AND the parser before any data is touched.

   Why "gate"? A gate lets only well-formed traffic through. Each
   stage raises the well-formedness bar one level: raw bytes -> tokens
   -> a validated Command. Whatever fails a bar never travels further.
\end{verbatim}

Gating is the pattern that makes layered systems safe. The lexer can
reject garbage cheaply, so the parser only ever sees sensible tokens;
the parser rejects wrong-shaped commands, so the storage engine only
ever sees validated operations. Each stage trusts its input, but the
trust is earned one gate back. MiniDB drives the whole pipeline this
way: read a line, lex it, parse it, route it, echo the reply, repeat.

\begin{lstlisting}[caption={The REPL pipeline (adapted from main.c).}]
#include <stdio.h>
#include <string.h>

/* Returns 1 if the line contains nothing but whitespace. */
static int is_blank_line(const char *line) {
    int i;

    for (i = 0; line[i] != '\0'; i++) {
        if (line[i] != ' ' && line[i] != '\t' &&
            line[i] != '\n' && line[i] != '\r') {
            return 0;
        }
    }
    return 1;
}

int main(void) {
    Token tokens[LEXER_MAX_TOKENS];   /* the caller owns the array */
    size_t token_count;
    Command cmd;
    char line[256];

    for (;;) {
        /* Prompt, and FLUSH: without this, a piped stdout would hold
           the prompt in a buffer while we wait for input it never
           shows. Interactive programs must flush.               */
        printf("minidb> ");
        fflush(stdout);

        /* EOF (Ctrl-D or a closed pipe) ends the session. */
        if (fgets(line, sizeof(line), stdin) == NULL) {
            break;
        }
        if (is_blank_line(line)) {
            continue;                  /* noise: skip silently */
        }

        /* Stage 1: LEX. A failed scan invalidates the stream. */
        token_count = 0;
        if (lex(line, tokens, &token_count) != LEX_OK) {
            printf("Lex error.\n");
            continue;
        }

        /* Stage 2: PARSE. The command struct is only filled on OK. */
        if (parse_command(tokens, token_count, &cmd) != PARSE_OK) {
            printf("Parse error. Type HELP for commands.\n");
            continue;
        }

        /* EXIT is REPL control flow, not a database operation. */
        if (cmd.type == COMMAND_EXIT) {
            printf("Goodbye!\n");
            break;
        }

        /* Stage 3: EXECUTE. MiniDB routes the Command to its B+ tree
           (Chapter 4) and prints the returned message. */
        execute_command(&cmd);
    }

    return 0;
}
\end{lstlisting}

\begin{verbatim}
   The loop, drawn as a state machine of its own:

        +----> print "minidb> " ; fflush
        |           |
        |           v
        |   fgets() -- EOF? ------> break
        |           |
        |           v
        |   blank? --------------> continue (skip)
        |           |
        |           v
        |   LEX -- error? --------> print, continue
        |           |
        |           v
        |   PARSE -- error? ------> print, continue
        |           |
        |           v
        |   EXIT? ----------------> break
        |           |
        |           v
        |   execute ; print reply  --> back to the top
        +-----------+
\end{verbatim}

The three details worth internalizing:

\begin{itemize}
  \item \emph{The stack owns the buffers.} \code{Token tokens[16]},
    \code{Command cmd}, and \code{char line[256]} are all locals. The
    lexer and parser write into caller-owned objects and allocate
    nothing. Memory reasoning in the front end is trivial: stack in,
    stack out, nothing to free.
  \item \emph{Every stage gates the next.} A lex error prints and
    skips parsing; a parse error prints and skips executing. A bad
    line can never reach the storage engine. The database is protected
    from the user by two walls of validation, and neither wall can be
    jumped.
  \item \emph{\code{EXIT} is REPL control flow, not a data command.}
    The REPL checks \code{cmd.type == COMMAND\_EXIT} before routing.
    Where the loop lives, the exit decision lives.
\end{itemize}

\section{The Complete Journey: ``INSERT 42 100''}

Let us follow one real line end to end, because the pipeline is the
point and the layers are where the meaning becomes structure.

\begin{verbatim}
   Step 1 -- the string.
   fgets() puts  'I','N','S','E','R','T',' ','4','2',' ','1','0','0',
                 '\n','\0'  into char line[256].

   Step 2 -- lex.
   'I' is a letter -> WORD mode: absorb INSERT, uppercase it, match
       the keyword table, emit { TOKEN_KEYWORD, KEYWORD_INSERT }.
   ' ' -> skipped.  '4' is a digit -> INTEGER mode: value = 4*10+2 = 42.
   ' ' -> skipped.  '1' is a digit -> INTEGER mode: value = 100.
   '\n' -> skipped; '\0' ends the loop; the END sentinel is appended.

   index:   0          1          2           3
          (KEYWORD)   (INTEGER)   (INTEGER)   (END)
           INSERT        42          100
   token_count == 4

   Step 3 -- parse.
   tokens[0] is a KEYWORD named INSERT -> arity 2:
   consume 42 and 100, require END, fill the struct:
       Command { type = COMMAND_INSERT, key = 42, value = 100 }

   Step 4 -- execute.
   dispatch on COMMAND_INSERT: the pair (42, 100) is placed in the
   storage engine -- in MiniDB, a B+ tree (our Chapter 4 structure).

   Step 5 -- reply.
   "inserted key 42 = 100" travels back to the REPL and is printed.

   One typed line, five stages, three structs, zero heap allocations
   in the front end, and the data lives in a balanced sorted tree.
\end{verbatim}

\begin{verbatim}
   A real session:

     minidb> INSERT 42 100
     inserted key 42 = 100
     minidb> SELECT 42
     key 42 = 100
     minidb> RANGE 1 50
     2 keys in [1, 50]
     minidb> DELETE 7
     deleted key 7
     minidb> SELECT 7
     key 7 not found
     minidb> EXIT
     Goodbye!
\end{verbatim}

\section{Layering: The Tree Never Learns the Language}

The deepest lesson of MiniDB is not any single function; it is the
\emph{direction of the dependencies}. The lexer knows nothing about
the parser. The parser knows nothing about the storage engine. The
tree stores \code{(int, int)} pairs and cannot tell \code{SELECT} from
\code{DELETE}---and does not need to. Each layer only understands the
value handed to it:

\begin{verbatim}
   The include graph is a strict, acyclic stack:

        main.c  (the REPL)
           |
           v
      database.h   <-- the seam: Command -> tree call -> reply
        /      \
       v        v
   parser.h   btree.h        (btree = our Chapter 4 B+ tree)
       |
       v
    lexer.h

   Read the arrows as "depends on":
     - lexer depends on nothing; it could feed any grammar.
     - parser validates streams; it has never heard of leaves.
     - tree stores pairs; it has never heard of tokens.

   Consequence, the sentence worth keeping: GROW THE LANGUAGE, AND THE
   DATABASE STAYS PUT. Add a new command and the tree does not change;
   rebalance the tree and the grammar does not change.
\end{verbatim}

That decoupling is the same front-end/back-end seam every real database
uses---the place where a ``planner'' would grow---and it is the payoff
of doing the parsing work before any storage work. The bytes become
tokens, tokens become a \code{Command}, a \code{Command} becomes tree
mutation, tree mutation becomes a reply. At each boundary, the
representation is \emph{more structured and less ambiguous}.

\begin{researchnote}[The theory of the state machine]
The lexer you just read is a \emph{finite automaton}, and the theory
behind recognizing text with one was laid down in a classic 1968 paper
by Ken Thompson: ``Regular Expression Search Algorithm.'' Thompson
showed how to compile a regular expression into a small machine that
walks the text one character at a time---exactly the discipline of
MiniDB's \code{lex()}.

\begin{verbatim}
   Thompson's insight, condensed:

       "A pattern can be compiled into a finite automaton that scans
        the input character by character, keeping only the state it
        is in; the state is all the history the machine needs."

   That is precisely MiniDB's lexer: four states, one current
   character, no lookahead, no backtracking. "The state is all the
   history" is why the lexer is a loop and not a pile of ifs -- and
   why it can reject bad input without ever having seen the future.
\end{verbatim}

Thompson's algorithm became the ancestor of every regular expression
engine and every lexer generator (lex, flex). Reading it is the
fastest way to understand why the state machine is the right shape
for character-level scanning---it is the shape with a sixty-year
pedigree.

\paperref{K.~Thompson, ``Regular Expression Search Algorithm,''
Communications of the ACM 11(6), 1968}{https://dl.acm.org/doi/10.1145/363347.363387}
\end{researchnote}

\section{Structures in Bytes}

The parser's structures are tiny, and their byte sizes are part of the
design---they are what make the fixed arrays safe.

\begin{verbatim}
   struct Token (lexer output)     12 bytes   offset size field
                                              0      4    type
                                              4      4    keyword
                                              8      4    int_value

   struct Command (parser output)  16 bytes   offset size field
                                              0      4    type
                                              4      4    key
                                              8      4    value
                                              12     4    key_end

   the REPL's buffers, in bytes:
     Token  tokens[16]     ->  16 * 12 = 192 bytes
     char   line[256]      ->  256 bytes
     Command cmd           ->  16 bytes

   LexResult and ParseResult enums: 4 bytes each (an int).
\end{verbatim}

Two consequences fall out of the sizes. First, a 256-byte line can
never overflow the 16-token array---each token needs at least one
byte, and 16 tokens is the enforced ceiling. Second, because every
field is a plain 4-byte int, a \code{Command} can be copied, compared,
and passed around byte-for-byte---exactly what the REPL does when it
routes it to the storage engine.

\section{Summary}

\begin{verbatim}
   The parsing toolbox, complete:

   token          the smallest meaningful unit
   lexer          characters -> tokens (a state machine)
   TOKEN_END      the sentinel that makes streams self-terminating
   parser         tokens -> Command (shape and arity validation)
   exact arity    consume N integers, then REQUIRE the sentinel
   proto-AST      a Command: one flat node for a flat grammar
   REPL           read, gate, execute, print -- one loop
   every stage gates the next   bad lines never reach storage

   Rules that keep parsers correct:
     1. draw the state machine before writing the loop;
     2. never read past a bound: sentinel, capacity check, buffer
        size all agree;
     3. overflow is a state transition to ERROR, not a silent wrap;
     4. on any lex/parse error, treat the stream as garbage;
     5. keep layers one-directional: the tree never learns the
        language.
\end{verbatim}

\begin{osconcept}
Parsing is where the untrusted bytes of the outside world first meet
your invariants, and the state machine is the pattern that makes the
meeting safe: every input, good or bad, ends in exactly one state. The
front-end/back-end seam---language on one side, storage on the
other---is the same seam every database draws, and it is the reason a
system can grow without its layers learning about each other.
\end{osconcept}

With a parser in hand we have the final major tool of the book. The
next chapter takes a step back from writing code to map the territory
every program we have built lives in: the Linux filesystem hierarchy---
what every directory is for, how the tree grew, and the operating
system concepts behind it.

